\section{Building Intuition: Accumulation and Random Mixing}
\label{sec:spin-intuition}

The outer gives sparse messages an initial weight floor; the inner must
turn this local protection into distance proportional to the full codeword
length. We compare two constructions with the same outer and permutation
to see what helps, and what obstructs, this spreading.

Choose an even block length $B$ and split the message into $M$ blocks of
$B/2$ bits each. The outer encodes block $i$ using a linear injection
$C_i:\F_2^{B/2}\to\F_2^B$. Setup samples these maps independently and
uniformly from all such injections. For each fixed nonzero block input,
its output is therefore uniform over the nonzero $B$-bit vectors.
Under the logarithmic block schedules below, with probability $1-o(1)$,
every nonzero block output has weight proportional to $B$.
In particular, this rules out weight-one outer outputs. But a message
supported on one block still produces at most $B$ nonzero bits, not a
constant fraction of the full length $N=MB$.

An independently sampled uniform permutation scatters the concatenated
outer outputs across the $N$ positions. Conditional on any fixed outer
word of weight $w$, every support of size $w$ is equally likely. The inner
then processes this scattered input. All setup choices are sampled once
and reused for every message.

\paragraph{The Minimal Inner: An Accumulator.}
For interleaved input $u$, set $y_0:=0$ and compute
\[
  y_i:=y_{i-1}+u_i,\qquad i=1,\ldots,N.
\]
Write $\Acc_N(u):=y$. With this inner, SPIN truncates
BAA~\cite{KolesnikovEtAl2026BA} to a single accumulator stage.
The inner has one bit of state. An input one starts a run of output ones,
but the next input one stops it. For example, input $00110000$ produces
output $00100000$: two adjacent ones leave only one nonzero output.
The outer's weight floor does not prevent this kind of cancellation.
An even-weight outer word of weight $w$ can be interleaved into adjacent
pairs, producing only $w/2$ output ones. More generally, each run of output ones starts at an odd-numbered input one
and ends just before the next input one. Unfavorable permutations place the
two input ones of each such pair close together,
leaving short runs of output ones separated by long runs of zeros.

The question is how often these unfavorable arrangements occur across
all nonzero messages. Counting the lengths of the output runs of ones and
zeros gives an exact input--output weight enumerator for the accumulator. Averaging
over the permutation and outer gives the expected number of low-weight
codewords. As the outer weight increases, a low-weight output requires
more short runs. Logarithmically growing outer blocks make these events
rare enough that the expected number of low-weight codewords tends to zero
at a positive relative distance.

Theorem~\ref{thm:independent-accumulator-spin} makes this concrete:
with its block schedule $B\ge3\log_2N$, the rate-$1/2$ family has relative
distance greater than $0.0037$ with probability $1-o(1)$ over setup.
Thus a one-bit inner suffices for linear distance, despite its easy
cancellation patterns.
With a terminated streaming outer, it also answers the question of
Bazzi, Mahdian, and Spielman~\cite{bms09} (Appendix~\ref{sec:bms-question}).

To obtain larger distance, one can permute and accumulate again, as in
BAA codes~\cite{KolesnikovEtAl2026BA}. The additional permutation breaks up
the first accumulator's runs, but incurs another pass of irregular data
movement. A stronger inner can instead sustain spreading after a single
permutation by making it harder to return to the zero state.

\paragraph{A Stronger Inner: Random Recursive Convolution.}
We keep the same outer and permutation, but use the recursive inner from
Expand--Convolute~\cite{RaghuramanRindalTanguy2023EC}. We derive a new,
exact expected input--output weight enumerator for this inner; it combines
directly with the outer spectrum in our proof framework.
Rather than storing one parity bit, the inner remembers the last $m$ outputs. Sample
independent $\alpha_i\gets\F_2^m$ at setup, independently of the outer and
interleaver, and set $y_j:=0$ for $j\le0$. Compute
\[
 y_i:=u_i+\langle\alpha_i,(y_{i-1},\ldots,y_{i-m})\rangle.
\]
For a fixed input, condition on the preceding outputs. Whenever the state
is nonzero, the independent coefficient vector makes the next output bit
uniform, regardless of the current input bit. A second input one therefore
does not simply end the spreading started by the first. Immediately after
an output one, returning to the zero state requires $m$ consecutive output
zeros. Each of these steps still has nonzero state, so this reset has
probability $2^{-m}$, provided $m$ positions remain. Once the state is zero,
the next input one makes it nonzero again.

This gives the proof two cases to control. A low-weight output either
spends substantial time in the zero state or contains unusually few ones
during its nonzero-state steps. Long memory makes resets rare, while the
outer and permutation supply scattered input ones that restart spreading.
Because each nonzero-state output bit is uniform, few ones during those
steps is also unlikely.
Tracking the number of trailing output zeros, up to $m$, and recording
input and output weights gives this enumerator. Summing over outer words
then bounds the probability that any nonzero message has low-weight output.
As for a random rate-$1/2$ code, the lower tails must be small enough to
overcome the $2^{N/2}-1$ nonzero messages.

Theorem~\ref{thm:independent-random-spin} combines this calculation with
logarithmic outer block size and memory to recover nearly the full
random-code distance. At rate $1/2$, it proves relative
distance greater than $0.11002$ with probability $1-o(1)$, close to the
Gilbert--Varshamov value of approximately $0.110028$. The complete enumerator and
the bounds for sparse and linear weights appear together in
Appendix~\ref{sec:random-spin}.

\paragraph{Toward Efficient Encoding.}
Random maps simplify the analysis but are costly to encode.
A dense random outer block takes $\Theta(B^2)$ bit operations;
over $N/B$ blocks with $B=\Theta(\log N)$, this is $\Theta(N\log N)$ work.
The accumulator itself is linear-time, whereas the random recursive inner
evaluates $m=\Theta(\log N)$ feedback taps at each position, adding another
$\Theta(N\log N)$ term. The uniform permutation also requires irregular
data movement.

Section~\ref{sec:structured-spin} replaces each costly component with a
structured alternative of linear work.
